% Section content template for individual Educative lessons
% This file contains the actual content for a specific section
% Generated from Educative API section content

% Section: Code for the Airline Management System
% Section ID: 5995883335516160
% Section Slug: code-for-the-airline-management-system
% Generated: 2025-12-14T13:18:05.354296

We've reviewed various aspects of the airline management system and observed the attributes associated with the problem using different UML diagrams. Let 's explore the more practical side of things, where we will work on implementing the airline management system using multiple languages. This is usually the last step in an object-oriented design interview process.
\\
We have chosen the following languages to write the skeleton code of the different classes present in the airline management system:

\begin{itemize}
\item
\protect\phantomsection\label{ZRALeHyL38ZrryY_nxrL6}
Java
\item
\protect\phantomsection\label{sZ1gzY1XnAsJ5C4RwoVqt}
C\#
\item
\protect\phantomsection\label{2dJx6DgHira96ALmK77iy}
Python
\item
\protect\phantomsection\label{JlTPQMu_-KFmNRrQhHInp}
C++
\item
\protect\phantomsection\label{9Zgx9WJ9E4Khrt-3Tbhak}
JavaScript
\end{itemize}
\\
If you want to skip directly to the implementation and see the complete workflow, simply scroll down or click~\hyperref[Executable-code-Airline-management-system]{Executable Code: Airline Management System}.

\subsection{The airline management system classes}\label{fxOaz7TAJwEnLE48rpQ8n}
\\
This section will provide the skeleton code of the classes designed in the class diagram lesson.
\begin{quote}
\textbf{Note:} For simplicity, we are not defining getter and setter functions. The reader can assume that all class attributes are private, accessed through their respective public getter methods, and modified only through their public methods.
\end{quote}
\subsubsection{Constants}\label{EIGA9GZVUoTiLDP8w036p}
\\
The following code defines the various enums and custom data types being used in the airline management system:
\begin{quote}
\textbf{Note:} JavaScript does not support enumerations, so we will use the\ Object.freeze()\ method as an alternative that freezes an object and prevents further modifications.
\end{quote}

\textbf{Constant definitions}

\begin{lstlisting}[language=Java]
public class Address {
  private int zipCode;
  private String streetAddress;
  private String city;
  private String state;
  private String country;
}

enum AccountStatus {
  ACTIVE,
  DISABLED,
  CLOSED,
  BLOCKED
}

enum SeatStatus {
  AVAILABLE,
  BOOKED,
  CHANCE
}

enum SeatType {
  REGULAR,
  ACCESSIBLE,
  EMERGENCY_EXIT,
  EXTRA_LEG_ROOM
}

enum SeatClass {
  ECONOMY,
  ECONOMY_PLUS,
  BUSINESS,
  FIRST_CLASS
}

enum FlightStatus {
  ACTIVE,
  SCHEDULED,
  DELAYED,
  LANDED,
  DEPARTED,
  CANCELED,
  DIVERTED,
  UNKNOWN
}

enum ReservationStatus {
  REQUESTED,
  PENDING,
  CONFIRMED,
  CHECKED_IN,
  CANCELED
}

enum PaymentStatus {
  PENDING,
  COMPLETED,
  FAILED,
  DECLINED,
  CANCELED,
  REFUNDED
}
\end{lstlisting}

\subsubsection{Account and passenger}\label{duDYtAPD24GfxxwnqwP4_}
\\
The \texttt{Account} class refers to an account for any user, including administrators, crew members, front desk officers, and customers. The
\texttt{Passenger} class represents the passengers in the airline system. The definition of this class is given below:

\textbf{The Account and Passenger classes}

\begin{lstlisting}[language=Java]
public class Account {
  private AccountStatus status;
  private int accountId;
  private String username;  
  private String password;

  public boolean resetPassword();
}

public class Passenger {
  private int passengerId;
  private String name;
  private String gender;
  private Date dateOfBirth;  
  private String passportNumber;
}
\end{lstlisting}

\subsubsection{Person}\label{rZQU1btcCeTmM8Dh6w8ow}
\\
\texttt{Person} is an abstract class that represents the various people or actors that can interact with the system. There are four types of persons: \texttt{Admin}, \texttt{Crew}, \texttt{FrontDeskOfficer}, and \texttt{Customer}. The implementation of the mentioned classes is shown below:

\textbf{Person and its derived classes}

\begin{lstlisting}[language=Java]
public abstract class Person {
  private String name;
  private Address address;
  private String email;
  private String phone;
  private Account account;
}

public class Admin extends Person {
  public boolean addAircraft(Aircraft aircraft);
  public boolean addFlight(Flight flight);
  public boolean cancelFlight(Flight flight);
  public boolean assignCrew(Flight flight);
  public boolean blockUser(User user);
  public boolean unblockUser(User user);
}

public class Crew extends Person {
  public List<FlightInstance> viewSchedule();
}

public class FrontDeskOfficer extends Person {
  public List<Itinerary> viewItinerary();
  public boolean createItinerary();
  public boolean createReservation();
  public boolean assignSeat();
  public boolean makePayment();
}

public class Customer extends Person {
  private int customerId;

  public List<Itinerary> viewItinerary();
  public boolean createItinerary();
  public boolean createReservation();
  public boolean assignSeat();
  public boolean makePayment();
}
\end{lstlisting}

\subsubsection{Seat and flight seat}\label{OHgBg0HPEEy-U8QzVlSjR}
\\
The \texttt{Seat} and \texttt{FlightSeat} are used to keep track of the customer's seat. Here , \texttt{Seat} is the physical seat in the aircraft, and \texttt{FlightSeat} is the seat assigned to a specific flight instance. The definition of these two classes is given below:

\textbf{The Seat and FlightSeat classes}

\begin{lstlisting}[language=Java]
public class Seat {
  private String seatNumber;
  private SeatType type;
  private SeatClass _class;
}

public class FlightSeat extends Seat {
  private double fare;
  private SeatStatus status;
  private String reservationNumber;
}
\end{lstlisting}

\subsubsection{Flight and flight instance}\label{uu_kBM_M2lr379s1c1kMX}
\\
The \texttt{Flight} and \texttt{FlightInstance} classes provide the details of the flight and its instances to the customer. The definition of these classes is given below:

\textbf{The Flight and FlightInstance classes}

\begin{lstlisting}[language=Java]
public class Flight {
  private String flightNo;
  private int durationMin;
  private Airport departure;
  private Airport arrival;
  private List<FlightInstance> instances;
}

public class FlightInstance {
  private Flight flight;
  private Date departureTime;
  private String gate;
  private FlightStatus status;
  private Aircraft aircraft;
  private List<FlightSeat> seats;
}
\end{lstlisting}

\subsubsection{Itinerary and flight reservation}\label{3EcaKLwBrmU3kxliGcOks}
\\
The \texttt{Itinerary} and \texttt{FlightReservation} classes are used to track itineraries and flights reserved by customers. Both classes are defined below:

\textbf{The Itinerary and FlightReservation classes}

\begin{lstlisting}[language=Java]
public class Itinerary {
  private Airport startingAirport;
  private Airport finalAirport;
  private Date creationDate;
  private List<FlightReservation> reservations;
  private List<Passenger> passengers;

  public boolean makeReservation();
  public boolean makePayment();
}

public class FlightReservation {
  private String reservationNumber;
  private FlightInstance flight;
  private HashMap<Passenger, FlightSeat> seatMap;
  private ReservationStatus status;
  private Date creationDate;

  public static FlightReservation fetchReservationDetails(String reservationNumber);
  public List<Passenger> getPassengers();
}
\end{lstlisting}

\subsubsection{Payment}\label{X526vRzSK5KV3s-gUqVJ3}
\\
The \texttt{Payment} class is another abstract class with two child classes: \texttt{Cash} and \texttt{CreditCard}. This utilizes the \texttt{PaymentStatus} enum to track the payment status. The definition of this class is provided below:

\textbf{Payment and its derived classes}

\begin{lstlisting}[language=Java]
public abstract class Payment {
  private int paymentId;
  private double amount;
  private PaymentStatus status;
  private Date timestamp;

  public abstract boolean makePayment();
}

public class Cash extends Payment {
    public boolean makePayment() {
        // functionality
    }
}

public class CreditCard extends Payment {
    private String nameOnCard;
    private String cardNumber;

    public boolean makePayment() {
        // functionality
    }
}
\end{lstlisting}

\subsubsection{Notification}\label{Bqb3JAkjP5oef7tysf05a}
\\
The \texttt{Notification} class is another abstract class responsible for sending notifications with two child classes: \texttt{SMSNotification} and \texttt{EmailNotification} . The implementation of this class is shown below:

\textbf{Notification and its derived classes }

\begin{lstlisting}[language=Java]
public abstract class Notification {
    private int notificationId;
    private Date createdOn;      
    private String content;

    public abstract void sendNotification(Account account);
}

class SmsNotification extends Notification {
    public void sendNotification(Account account) {
        // functionality 
    }
}

class EmailNotification extends Notification {
    public void sendNotification(Account account) {
        // functionality 
    }
}
\end{lstlisting}

\subsubsection{Search and catalog}\label{OELfoo90MCtsrHnuznL_n}
\\
The \texttt{SearchCatalog} class contains flight instance information and implements the \texttt{Search} interface to enable search functionality based on the specified criteria. Both classes are defined below:

\textbf{The Search interface and the SearchCatalog class}

\begin{lstlisting}[language=Java]
public interface Search {
  // Interface method (does not have a body)
  public List<FlightInstance> searchFlight(Airport source, Airport dest, Date arrival, Date departure);
}

public class SearchCatalog implements Search {
  private HashMap<Quartet<Airport, Airport, Date, Date>, List<FlightInstance>> flights;

  public List<FlightInstance> searchFlight(Airport source, Airport dest, Date arrival, Date departure) {
    // functionality
  }
}
\end{lstlisting}

\subsubsection{Airport, aircraft, and airline}\label{WHgbf8-XKpaHt_gUP8_kI}
\\
This section contains classes such as \texttt{Airport}, \texttt{Aircraft}, and \texttt{Airline} that comprise the infrastructure of our airline management system. Here , the \texttt{Airline} class is a Singleton. The definition of these classes is given below:

\textbf{The Airport, Aircraft, and Airline classes}

\begin{lstlisting}[language=Java]
public class Airport {
  private String name;
  private String code;
  private Address address;
  private List<Flight> flights;
}

public class Aircraft {
  private String name;
  private String code;
  private String model;
  private int seatCapacity;
  private List<Seat> seats;
}

public class Airline {
  private String name;
  private String code;
  private List<Flight> flights;
  private List<Aircraft> aircrafts;
  private List<Crew> crew;
  
  // The Airline is a singleton class that ensures it will have only one active instance at a time
  private static Airline airline = null;
  
  // Created a static method to access the singleton instance of Airline class
  public static Airline getInstance() {
    if (airline == null) {
      airline = new Airline();
    }
    return airline;
  }
}
\end{lstlisting}

\subsection{Executable code: Airline management system}\label{hgwDc5tfqfmWkL3eAyJr8}
\\
Below is a fully self-contained, runnable program in Java, C\#, C++, Python, and JavaScript that demonstrates the core workflows of the Airline Management System. The main driver code of the system resides in the~\texttt{Driver.java},~\texttt{Driver.cs},~\texttt{Driver.py},~\texttt{Driver.cpp}, and~\texttt{Driver.js}~for all respective languages. You can click the ``Run'' button to execute the codes.

\subsubsection{What does this code show}\label{0iIWSHM84aHEXUV_7iXLT}

\begin{itemize}
\item
\protect\phantomsection\label{lZSH5-mzHIs8g6dW6Lcxy}
\textbf{System initialization:} Initializes an admin, a customer, and a crew member, each with their own \texttt{Account}, \texttt{Address}, and role-specific objects. Two airports are set up, along with a Boeing 737 aircraft preloaded with seats. An initial flight is created between the two airports, and a flight instance is generated using the aircraft. This establishes the baseline for simulating operations such as booking, payments, and crew assignments.
\item
\protect\phantomsection\label{ikL53EhQVo2l1n5uNfqVQ}
\textbf{Scenario 1 (Customer booking and payment):} The customer books a seat on the initialized flight by creating a reservation and associating it with a passenger. A payment is made using a credit card, and an email notification is sent upon successful payment, demonstrating the reservation and transaction flow through user interaction.
\item
\protect\phantomsection\label{RPEljgTplN7Qth4nmAdyQ}
\textbf{Scenario 2 (Admin flight and crew management):} The admin creates a new aircraft and adds a new flight using it. A flight instance is created, and the admin assigns the existing crew member to this instance. This highlights how administrators manage operational logistics, such as flights and personnel assignments.
\item
\protect\phantomsection\label{kehlXwfqd2rakjLQr-3Ia}
\textbf{Scenario 3 (Crew schedule viewing):} The crew member views their schedule through the \texttt{viewSchedule()} method, which retrieves their assigned flights via the admin's internal crew scheduling system. This simulates how crew members access up-to-date information about their assignments.
\end{itemize}
\begin{quote}
\textbf{Hint: Try customizing the code!}\\
This code is designed not just for reading, but for experimentation. You can~edit, extend, or modify~any part of the example.
\end{quote}
\subsection{Code executable}\label{OxJPHl4bw81G2Wb_Kf7ph}

\textbf{Executable Code: Airline Management System}

\begin{lstlisting}[language=Java]
import java.util.*;

public class Driver {

    public static void main(String[] args) {
        System.out.println("== Initializing system... ==");
        
        // Create addresses
        Address address1 = new Address(12345, "Street 1", "CityA", "StateA", "CountryA");
        Address address2 = new Address(67890, "Street 2", "CityB", "StateB", "CountryB");
        System.out.println("Addresses created.");
        
        // Airports
        Airport airport1 = new Airport("Alpha Airport", "ALP", address1, new ArrayList<>());
        Airport airport2 = new Airport("Beta Airport", "BET", address2, new ArrayList<>());
        System.out.println("Airports created: " + airport1.getName() + " and " + airport2.getName());
        
        // Aircraft and seats
        List<Seat> aircraftSeats = new ArrayList<>();
        for (int i = 1; i <= 10; i++) {
            aircraftSeats.add(new Seat("A" + i, SeatType.REGULAR, SeatClass.ECONOMY));
        }
        Aircraft aircraft = new Aircraft("Boeing 737", "B737", "737-800", 10, aircraftSeats);
        Airline.getInstance().addAircraft(aircraft);
        System.out.println("Aircraft created and added: " + aircraft.getName());
        
        // Accounts
        Account customerAccount = new Account(101, "jane_doe", "password");
        Account adminAccount = new Account(201, "admin1", "adminpass");
        Account crewAccount = new Account(301, "pilot_guy", "flyhigh");
        System.out.println("User accounts created.");
        
        // People
        Customer customer = new Customer(1, "Jane Doe", address1, "jane@example.com", "111-222-3333", customerAccount);
        Admin admin = new Admin("Admin One", address2, "admin@example.com", "444-555-6666", adminAccount);
        Crew crew = new Crew("Captain Sky", address1, "pilot@example.com", "777-888-9999", crewAccount);
        Airline.getInstance().addCrew(crew);
        System.out.println("Customer, Admin, and Crew created.");
        
        // Flight and flight instance
        Flight flight = new Flight("FL123", 120, airport1, airport2, new ArrayList<>());
        FlightInstance flightInstance = new FlightInstance(flight, new Date(), "G5", FlightStatus.SCHEDULED, aircraft, new ArrayList<>());
        
        for (Seat seat : aircraftSeats) {
            flightInstance.getSeats().add(new FlightSeat(seat.getSeatNumber(), seat.getType(), seat.getSeatClass(), 150.0));
        }
        flight.getInstances().add(flightInstance);
        admin.addFlight(flight);
        System.out.println("Initial flight created and flight instance added.");
        
        System.out.println("System initialized.
");


        // ========= Scenario 1 =========
        System.out.println("== Scenario 1: Customer books and pays ==");

        Passenger p1 = new Passenger(1, "Jane Doe", "Female", createDate(1990, 5, 5), "P12345");
        FlightSeat selectedSeat = flightInstance.getSeats().get(0);
        selectedSeat.setStatus(SeatStatus.BOOKED);
        selectedSeat.setReservationNumber("RESV001");

        HashMap<Passenger, FlightSeat> seatMap = new HashMap<>();
        seatMap.put(p1, selectedSeat);

        FlightReservation reservation = new FlightReservation("RESV001", flightInstance, seatMap, ReservationStatus.CONFIRMED, new Date());

        List<FlightReservation> reservations = new ArrayList<>();
        reservations.add(reservation);
        List<Passenger> passengers = new ArrayList<>();
        passengers.add(p1);

        Itinerary itinerary = new Itinerary(airport1, airport2, new Date(), reservations, passengers);
        itinerary.makeReservation();

        CreditCard card = new CreditCard(1, selectedSeat.getFare(), "Jane Doe", "4111111111111111");
        customer.makePayment(card);

        System.out.println();

        // ========= Scenario 2 =========
        System.out.println("== Scenario 2: Admin adds a new flight and assigns crew ==");

        Aircraft newAircraft = new Aircraft("Airbus A320", "A320", "A320-200", 5, new ArrayList<>());
        admin.addAircraft(newAircraft);

        Flight newFlight = new Flight("FL456", 180, airport2, airport1, new ArrayList<>());
        FlightInstance newInstance = new FlightInstance(newFlight, new Date(), "H2", FlightStatus.SCHEDULED, newAircraft, new ArrayList<>());
        newFlight.getInstances().add(newInstance);

        admin.addFlight(newFlight);
        admin.assignCrew(crew, newInstance);

        System.out.println();

        // ========= Scenario 3 =========
        System.out.println("== Scenario 3: Crew views schedule ==");

        List<FlightInstance> crewSchedule = crew.viewSchedule();
        for (FlightInstance fi : crewSchedule) {
            System.out.println("Assigned to Flight " + fi.getFlight().getFlightNo() + " at Gate " + fi.getGate());
        }
    }

    // Utility method to safely create a Date from year, month, day
    private static Date createDate(int year, int month, int day) {
        Calendar cal = Calendar.getInstance();
        cal.set(year, month - 1, day); // Java months are 0-based
        return cal.getTime();
    }
}
\end{lstlisting}

\subsection{Wrapping up}\label{fQEqvXHpCh0-mP6Zw64no}
\\
We've explored the complete design of the airline management system in this chapter. We 've looked at how a basic airline management system can be visualized using various UML diagrams and designed using object-oriented principles and design patterns.


% End of section content